%============================================================
% CH1 · LESSON 5 — IN PLACE VERSUS EXTRA MEMORY
%============================================================
\subsection{In Place versus Extra Memory}
\label{ssec:ch1-inplace}

Many array tasks produce a \emph{new} arrangement of the elements: remove some, move some to the end, reverse, rotate, replace each value by a function of its neighbours. There are two ways to produce it:
\begin{itemize}[leftmargin=*]
    \item \textbf{with extra memory:} build the result in a second array, then copy it back (or return it). Simple, hard to get wrong, $\bigO{n}$ extra memory.
    \item \textbf{in place:} rearrange the input array itself using only $\bigO{1}$ extra variables. Required when the statement says so (``do not allocate another array'', ``modify \texttt{nums} in place with $\bigO{1}$ extra memory''), and useful when memory is tight.
\end{itemize}
The danger of working in place is always the same: \textbf{a write destroys a value that is still needed}. This lesson collects the four standard ways to avoid that --- swapping, reversal with two pointers, a write pointer that stays behind the read pointer, and packing two values into one cell --- and ends with the property that distinguishes good in-place algorithms from bad ones: stability.

%------------------------------------------------------------
\subsubsection*{1. Why ``in place'' is harder: the overwrite problem}

Task: shift every element one position to the right (the last one falls off), $a_i \gets a_{i-1}$. The obvious left-to-right loop is wrong:

\begin{diagrambox}
\begin{tikzpicture}
    \begin{scope}
        \dparray{x}{0}{wrong}{5/acell, 5/awrite, 5/awrite, 5/awrite, 5/awrite}
        \node[note, anchor=west, text width=4.6cm, align=left] at (4.8,0) {left to right: \texttt{a[1] = a[0]} destroys the old \texttt{a[1]} before \texttt{a[2]} reads it --- everything becomes 5};
    \end{scope}
    \begin{scope}[yshift=-1.3cm]
        \dparray{y}{0}{right}{5/acell, 5/awrite, 8/awrite, 1/awrite, 6/awrite}
        \node[note, anchor=west, text width=4.6cm, align=left] at (4.8,0) {right to left: \texttt{a[4] = a[3]} first, then \texttt{a[3] = a[2]}, \ldots{} each value is read before it is overwritten};
    \end{scope}
    \node[font=\scriptsize\color{darkgray}] at (2.25,0.75) {input $(5, 8, 1, 6, 2)$};
\end{tikzpicture}
\figcap{The direction of an in-place loop decides whether it reads old or new values. Moving data right, scan from the right; moving data left, scan from the left.}
\end{diagrambox}

\begin{insightbox}[The one question to ask]
Before every in-place write \texttt{a[w] = \ldots}, ask: \emph{is the old value of \texttt{a[w]} still needed by a later step?} If yes, either save it first (swap), process in the other direction, or keep it inside the cell in encoded form (Section~5).
\end{insightbox}

%------------------------------------------------------------
\subsubsection*{2. Swap}

\texttt{swap(a[i], a[j])} exchanges two cells through one temporary variable, so neither value is lost:
\begin{lstlisting}[style=snippet]
int tmp = a[i]; a[i] = a[j]; a[j] = tmp;     // what std::swap(a[i], a[j]) does
\end{lstlisting}
Every in-place permutation of an array (reversal, rotation, partition, sorting) can be written as a sequence of swaps.

\begin{warningbox}[The XOR swap]
\texttt{a[i] \^{}= a[j]; a[j] \^{}= a[i]; a[i] \^{}= a[j];} swaps without a temporary --- and sets the cell to $0$ when $i = j$, because then all three lines operate on the same cell ($x \oplus x = 0$). It is also not faster. Use \texttt{std::swap}.
\end{warningbox}

%------------------------------------------------------------
\subsubsection*{3. Reversal with two pointers}

To reverse $[l, r)$, swap the outermost pair, then move both pointers inwards:
\begin{lstlisting}[style=snippet]
for (int i = l, j = r - 1; i < j; i++, j--) swap(a[i], a[j]);    // or std::reverse(a + l, a + r)
\end{lstlisting}

\begin{diagrambox}
\begin{tikzpicture}
    \begin{scope}
        \dparray{r}{0}{start}{1/aread, 2/acell, 3/acell, 4/acell, 5/acell, 6/acell, 7/aread}
        \aptr{r0}{i}{successgreen} \aptr{r6}{j}{accentorange}
    \end{scope}
    \begin{scope}[yshift=-1.75cm]
        \dparray{s}{0}{2 swaps}{7/adone, 6/adone, 3/aread, 4/acell, 5/aread, 2/adone, 1/adone}
        \aptr{s2}{i}{successgreen} \aptr{s4}{j}{accentorange}
    \end{scope}
    \begin{scope}[yshift=-3.5cm]
        \dparray{t}{0}{done}{7/adone, 6/adone, 5/adone, 4/adone, 3/adone, 2/adone, 1/adone}
        \aptr{t3}{i = j}{successgreen}
    \end{scope}
    \node[note, anchor=west, text width=3.8cm, align=left] at (6.6,-1.75)
        {invariant: outside $[i, j]$ the array is already reversed; the loop stops when $i \ge j$ after $\lfloor (r-l)/2 \rfloor$ swaps};
\end{tikzpicture}
\figcap{Reversing $(1,\dots,7)$: three swaps; the middle element stays where it is.}
\end{diagrambox}

\begin{invariantbox}[Why the reversal is correct]
Reversing $[l, r)$ means: the element at position $p$ moves to $l + r - 1 - p$. \textbf{Invariant:} before each check of $i < j$, we have $i + j = l + r - 1$, every position $p < i$ or $p > j$ holds the original element of position $l + r - 1 - p$, and $[i, j]$ is untouched. A swap of $a_i$ and $a_j$ puts both into their final positions (since $j = l + r - 1 - i$), and moving $i \gets i + 1$, $j \gets j - 1$ keeps $i + j$ constant. When $i \ge j$, at most the middle cell is left and it is its own mirror image. \hfill$\blacksquare$
\end{invariantbox}

Reversal is the building block of rotation: rotating right by $k$ is ``reverse all, reverse the first $k$, reverse the rest'' --- Pattern~3, Section~\ref{sec:ch1-p3}.

%------------------------------------------------------------
\subsubsection*{4. Read and write pointers}

To \emph{remove} elements in place --- keep only those that satisfy a condition, in their original order --- use two indices moving in the same direction: a \textbf{read pointer} $r$ that visits every element, and a \textbf{write pointer} $w$ that marks where the next kept element goes.
\begin{lstlisting}[style=snippet]
int w = 0;
for (int r = 0; r < n; r++)
    if (keep(a[r])) a[w++] = a[r];    // a[0..w) = kept elements of a[0..r), in order
// new length is w
\end{lstlisting}

\begin{diagrambox}
\begin{tikzpicture}
    \dpindex{wi}{1.35}{}{0}{7}
    \dparray{w}{0.6}{\texttt{a}}{4/akeep, 2/akeep, 6/akeep, 7/aout, 3/aout, 8/aread, 1/acell, 5/acell}
    \aptr{w3}{w}{accentorange}
    \aptr{w5}{r}{successgreen}
    \arange{w0}{w2}{kept so far}{successgreen}
    \node[note, anchor=west, text width=3.8cm, align=left] at (7.5,0.6)
        {keep the even values; cells $[w, r)$ are free to be overwritten because they have already been read};
\end{tikzpicture}
\figcap{Read/write pointers: $w \le r$ always, so a write can only hit a cell that was already read.}
\end{diagrambox}

The whole safety argument is the inequality $w \le r$: the write position never overtakes the read position, so no unread element is ever destroyed. This is Pattern~2 (Section~\ref{sec:ch1-p2}), and Problems~\ref{ssec:ch1-move-zeroes} and~\ref{ssec:ch1-remove-duplicates-ii} are two applications. The mirror image --- \emph{inserting} or duplicating elements in place --- needs the write pointer to run \emph{ahead}, so it is done from the back, right to left (Exercise~E2.1).

%------------------------------------------------------------
\subsubsection*{5. Encoding two values in one cell}

Sometimes every new value depends on \emph{old} values of other cells, in no convenient direction. Example: replace every element by the maximum of itself and its two neighbours in the \emph{original} array,
\[
b_i = \max(a_{i-1}, a_i, a_{i+1}) \qquad (\text{missing neighbours ignored}).
\]
Computing $b_i$ needs the old $a_{i-1}$, which a left-to-right loop has just overwritten; a right-to-left loop has the same problem with $a_{i+1}$. With extra memory it is trivial. In place, if the values are \textbf{small non-negative integers}, we can keep the old and the new value in the same cell:
\[
\text{cell} = \text{old} + M \cdot \text{new}, \qquad 0 \le \text{old}, \text{new} < M
\quad\Longrightarrow\quad
\text{old} = \text{cell} \bmod M, \quad \text{new} = \lfloor \text{cell} / M \rfloor .
\]
Adding a multiple of $M$ does not change the remainder modulo $M$, so a cell that has already been updated still reveals its old value. After the pass, a second pass divides every cell by $M$.

\begin{diagrambox}
\begin{tikzpicture}
    \dpindex{ei}{1.35}{$i$}{0}{7}
    \dparray{e}{0.6}{old}{3/acell, 1/acell, 4/acell, 1/acell, 5/acell, 9/acell, 2/acell, 6/acell}
    \dparray{f}{-0.25}{cell}{33/awrite, 41/awrite, 44/awrite, 51/aread, 5/acell, 9/acell, 2/acell, 6/acell}
    \dparray{g}{-1.1}{final}{3/adone, 4/adone, 4/adone, 5/adone, 9/adone, 9/adone, 9/adone, 6/adone}
    \node[note, anchor=west, text width=3.6cm, align=left] at (7.5,-0.25)
        {$M = 10$: after step $i = 3$, cell 2 holds $44 = 4 + 10\cdot 4$; its old value $44 \bmod 10 = 4$ is still readable};
\end{tikzpicture}
\figcap{Neighbour maximum of $(3,1,4,1,5,9,2,6)$ in place, with $M = \max + 1 = 10$. Blue cells are encoded; the final row is \texttt{cell}$/M$.}
\end{diagrambox}

\begin{lstlisting}[style=cpp]
#include <bits/stdc++.h>
using namespace std;

void reverseRange(vector<int>& a, int l, int r) {   // reverses a[l..r)
    for (int i = l, j = r - 1; i < j; i++, j--) swap(a[i], a[j]);
}

// b[i] = max(a[i-1], a[i], a[i+1]) of the ORIGINAL array, in place.
// Needs 0 <= a[i] < M; cell i holds old + M * new while both are needed.
void neighbourMax(vector<int>& a) {
    int n = a.size();
    int M = *max_element(a.begin(), a.end()) + 1;
    for (int i = 0; i < n; i++) {
        int best = a[i] % M;                                 // old value of a[i]
        if (i > 0) best = max(best, a[i - 1] % M);           // old value, even if a[i-1] is encoded
        if (i + 1 < n) best = max(best, a[i + 1] % M);
        a[i] += M * best;                                    // fits in int if M*M < 2^31
    }
    for (int i = 0; i < n; i++) a[i] /= M;                  // keep only the new values
}

int main() {
    vector<int> a = {1, 2, 3, 4, 5, 6, 7};
    reverseRange(a, 0, 7);
    reverseRange(a, 2, 5);
    for (int x : a) cout << x << " ";        // 7 6 3 4 5 2 1
    cout << "\n";
    vector<int> b = {3, 1, 4, 1, 5, 9, 2, 6};
    neighbourMax(b);
    for (int x : b) cout << x << " ";        // 3 4 4 5 9 9 9 6
    cout << "\n";
}
\end{lstlisting}

Conditions and alternatives:
\begin{itemize}[leftmargin=*]
    \item The encoding needs $\text{old} + M \cdot \text{new} < 2^{31}$, i.e.\ roughly $M < 46341$ for \texttt{int}; use \texttt{long long} cells or another trick for larger values.
    \item Typical $M$: $n$ when the values are indices in $[0, n)$ --- this is how Problem~\ref{ssec:ch1-build-permutation} of Pattern~4 (Section~\ref{sec:ch1-p4}) builds \texttt{a[a[i]]} in place.
    \item When the values are \emph{positive}, the sign bit is a free one-bit flag: negate $a_v$ to mark ``value $v$ was seen'' and use $|a_v|$ for the original value (Chapter~\ref{ch:index_tricks}).
    \item When all values are bounded by a small number of bits, the high bits of the cell can hold the second value: \texttt{cell = old | (new << 16)}.
\end{itemize}

%------------------------------------------------------------
\subsubsection*{6. Stability}

\begin{definitionbox}[Stable rearrangement]
A rearrangement is \textbf{stable} if elements that the rule does not distinguish keep their original relative order. ``Even values first'' is stable if the even values appear in their original order and so do the odd ones.
\end{definitionbox}

Stability matters whenever elements carry more than the key you rearrange by: records sorted by one field and then by another, or a problem that says ``keep the relative order of the non-zero elements'' (Problem~\ref{ssec:ch1-move-zeroes}). Two in-place-looking methods for ``evens first'' behave differently:

\begin{lstlisting}[style=cpp]
#include <bits/stdc++.h>
using namespace std;

// Put the even values first. Labels a,b,c,... show where each element came from.
int main() {
    vector<pair<int,char>> v = {{4,'a'}, {7,'b'}, {2,'c'}, {9,'d'}, {6,'e'}, {3,'f'}, {8,'g'}};
    // (1) read/write compaction into a copy: stable
    vector<pair<int,char>> st;
    for (auto& p : v) if (p.first % 2 == 0) st.push_back(p);
    for (auto& p : v) if (p.first % 2 != 0) st.push_back(p);
    // (2) swap from both ends: in place, NOT stable
    vector<pair<int,char>> un = v;
    int i = 0, j = (int)un.size() - 1;
    while (i < j) {
        if (un[i].first % 2 == 0) i++;
        else if (un[j].first % 2 != 0) j--;
        else swap(un[i++], un[j--]);
    }
    for (auto& p : st) cout << p.first << p.second << " ";   // 4a 2c 6e 8g 7b 9d 3f
    cout << "\n";
    for (auto& p : un) cout << p.first << p.second << " ";   // 4a 8g 2c 6e 9d 3f 7b
    cout << "\n";
}
\end{lstlisting}

\begin{diagrambox}
\begin{tikzpicture}
    \dparray{s0}{0.9}{input}{4a/acell, 7b/aout, 2c/acell, 9d/aout, 6e/acell, 3f/aout, 8g/acell}
    \dparray{s1}{0}{stable}{4a/akeep, 2c/akeep, 6e/akeep, 8g/akeep, 7b/aout, 9d/aout, 3f/aout}
    \dparray{s2}{-0.9}{two ends}{4a/akeep, 8g/aread, 2c/akeep, 6e/akeep, 9d/aout, 3f/aout, 7b/aread}
    \node[note, anchor=west, text width=3.6cm, align=left] at (6.6,0)
        {orange: elements that left their relative order --- \texttt{8g} jumped over \texttt{2c} and \texttt{6e}, \texttt{7b} ended up last};
\end{tikzpicture}
\figcap{Both outputs put the evens first; only the copy-based one keeps $a, c, e, g$ and $b, d, f$ in order.}
\end{diagrambox}

\begin{table}[H]
\centering\small
\renewcommand{\arraystretch}{1.2}
\begin{tabular}{@{}lccc@{}}
\toprule
\textbf{method} & \textbf{extra memory} & \textbf{stable} & \textbf{where} \\
\midrule
copy into a new array & $\bigO{n}$ & yes & always available \\
read/write pointers (remove or compact) & $\bigO{1}$ & yes (for the kept part) & Pattern~2 \\
swaps from both ends (partition) & $\bigO{1}$ & no & Chapter~\ref{ch:partitioning} \\
\texttt{std::stable\_partition}, \texttt{std::stable\_sort} & $\bigO{n}$ buffer & yes & library \\
\bottomrule
\end{tabular}
\end{table}

\begin{warningbox}[\texttt{vector::erase} in a loop]
\begin{lstlisting}[style=snippet]
for (int i = 0; i < (int)v.size(); )              // WRONG idea: O(n^2)
    if (v[i] == 0) v.erase(v.begin() + i); else i++;
\end{lstlisting}
Each \texttt{erase} shifts the whole tail left: up to $n$ erases of $\bigO{n}$ each. For $n = 10^5$ zeros that is $5\cdot10^9$ moves. The read/write loop does the same job in one pass; the library spells it \texttt{v.erase(remove(v.begin(), v.end(), 0), v.end())}.
\end{warningbox}

\begin{checklistbox}[Lesson summary]
\begin{itemize}[leftmargin=*]
    \item[$\square$] In place = $\bigO{1}$ extra memory. The bug to avoid: overwriting a value that is still needed.
    \item[$\square$] Choose the loop direction so that every cell is read before it is written (data moves right $\Rightarrow$ scan from the right).
    \item[$\square$] Swap and two-pointer reversal permute without loss; read/write pointers remove without loss because $w \le r$.
    \item[$\square$] Small non-negative values: store $\text{old} + M\cdot\text{new}$, read \texttt{\% M} and \texttt{/ M}.
    \item[$\square$] Ask whether the order of equal (or kept) elements must be preserved; swap-from-both-ends is not stable.
\end{itemize}
\end{checklistbox}
